% GENERATED FILE. Edit canonical lessons, then run npm run generate:books.
% canonical-source-hash: fnv1a64:a31fbee01cf44e9b
% canonical-lessons: AR-C431-nafaq, AR-C431-ishara, AR-C431-murur, AR-C431-izdiham, AR-C431-inatafa

\chapter{Traffic}
\label{ch:ar-traffic}
\hlchaptermodality{431}

\begin{chapteropening}
\textbf{By the end of this chapter:} \emph{I can say a tunnel, a signal, traffic, congestion and to turn off.}

Complete the last lesson of chapter 431: I can say a tunnel, a signal, traffic, congestion and to turn off.
\end{chapteropening}

\section[\texorpdfstring{\ar{نفق} (nafaq)}{نفق (nafaq)}]{\ar{نفق} (nafaq) — a tunnel}

\label{lesson:AR-C431-nafaq}



\begin{quote}
Before the new one: say the Arabic for a dinghy, then the Arabic for to row.
\end{quote}

\subsection*{You'll want to know: \ar{نفق}}
\textbf{\ar{نفق}} — \emph{nafaq} — \textquotedblleft{}a tunnel\textquotedblright{}.

\begin{quote}
\ar{دخل} \ar{القطار} \ar{نفقا} \ar{طويلا} \ar{تحت} \ar{الجبل}.
\end{quote}

\begin{quote}
\emph{dakhala al-qiṭār nafaqan ṭawīlan taḥta al-jabal.}
\end{quote}

\begin{quote}
The train went into a long tunnel under the mountain.
\end{quote}

Plural: \ar{أنفاق}, \emph{anfāq}.

\begin{cousinweb}
\textbf{\ar{نفق}} (\emph{nafaq}) is built on the native Arabic root \emph{n-f-q}, \textquotedblleft{}to pass away, to run out\textquotedblright{} (Lane). Lane gives \emph{nafaq} as a passage under the ground.
\end{cousinweb}

\begin{grammarlens}[title={What you've built}]
One more word for getting around.
\end{grammarlens}

\subsection*{Your turn}
\begin{itemize}
\raggedright
  \item \emph{Say it:} \emph{nafaq}
  \item \emph{Say it:} \emph{nafaq}, once more
  \item \emph{Say it:} \emph{dakhala al-qiṭār nafaqan ṭawīlan taḥta al-jabal.} — The train went into a long tunnel under the mountain.
  \item \emph{Recall:} say the Arabic for a dinghy, then the Arabic for to row, then say \emph{nafaq} again
\end{itemize}

\begin{tcolorbox}[breakable,colback=teal!4,colframe=teal!35!black,title={Before you move on}]
What does \ar{نفق} mean? (\textquotedblleft{}A tunnel\textquotedblright{}.) Say it once more.
\end{tcolorbox}
\section[\texorpdfstring{\ar{إشارة} (ishāra)}{إشارة (ishāra)}]{\ar{إشارة} (ishāra) — a signal; traffic lights}

\label{lesson:AR-C431-ishara}



\begin{quote}
Before the new one: say the Arabic for to row, then the Arabic for a tunnel.
\end{quote}

\subsection*{You'll want to know: \ar{إشارة}}
\textbf{\ar{إشارة}} — \emph{ishāra} — \textquotedblleft{}a signal; traffic lights\textquotedblright{}.

\begin{quote}
\ar{وقفت} \ar{السيارة} \ar{عند} \ar{الإشارة} \ar{الحمراء}.
\end{quote}

\begin{quote}
\emph{waqafat as-sayyāra ʿinda al-ishāra al-ḥamrāʾ.}
\end{quote}

\begin{quote}
The car stopped at the red light.
\end{quote}

\begin{cousinweb}
\textbf{\ar{إشارة}} (\emph{ishāra}) is built on the native Arabic root \emph{sh-w-r}, \textquotedblleft{}to point, to make a sign\textquotedblright{} (Lane). It is the verbal noun of Form IV \emph{ashāra}.
\end{cousinweb}

\begin{grammarlens}[title={What you've built}]
One more word for getting around.
\end{grammarlens}

\subsection*{Your turn}
\begin{itemize}
\raggedright
  \item \emph{Say it:} \emph{ishāra}
  \item \emph{Say it:} \emph{ishāra}, once more
  \item \emph{Say it:} \emph{waqafat as-sayyāra ʿinda al-ishāra al-ḥamrāʾ.} — The car stopped at the red light.
  \item \emph{Recall:} say the Arabic for to row, then the Arabic for a tunnel, then say \emph{ishāra} again
\end{itemize}

\begin{tcolorbox}[breakable,colback=teal!4,colframe=teal!35!black,title={Before you move on}]
What does \ar{إشارة} mean? (\textquotedblleft{}A signal; traffic lights\textquotedblright{}.) Say it once more.
\end{tcolorbox}
\section[\texorpdfstring{\ar{مرور} (murūr)}{مرور (murūr)}]{\ar{مرور} (murūr) — traffic; passing}

\label{lesson:AR-C431-murur}



\begin{quote}
Before the new one: say the Arabic for a tunnel, then the Arabic for a signal.
\end{quote}

\subsection*{You'll want to know: \ar{مرور}}
\textbf{\ar{مرور}} — \emph{murūr} — \textquotedblleft{}traffic; passing\textquotedblright{}.

\begin{quote}
\ar{كان} \ar{المرور} \ar{خفيفا} \ar{يوم} \ar{الجمعة}.
\end{quote}

\begin{quote}
\emph{kāna al-murūr khafīfan yawm al-jumʿa.}
\end{quote}

\begin{quote}
Traffic was light on Friday.
\end{quote}

\begin{cousinweb}
\textbf{\ar{مرور}} (\emph{murūr}) is built on the native Arabic root \emph{m-r-r}, \textquotedblleft{}to pass\textquotedblright{} (Lane). It is the verbal noun of \emph{marra}, \textquotedblleft{}to pass by\textquotedblright{}.
\end{cousinweb}

\begin{grammarlens}[title={What you've built}]
One more word for getting around.
\end{grammarlens}

\subsection*{Your turn}
\begin{itemize}
\raggedright
  \item \emph{Say it:} \emph{murūr}
  \item \emph{Say it:} \emph{murūr}, once more
  \item \emph{Say it:} \emph{kāna al-murūr khafīfan yawm al-jumʿa.} — Traffic was light on Friday.
  \item \emph{Recall:} say the Arabic for a tunnel, then the Arabic for a signal, then say \emph{murūr} again
\end{itemize}

\begin{tcolorbox}[breakable,colback=teal!4,colframe=teal!35!black,title={Before you move on}]
What does \ar{مرور} mean? (\textquotedblleft{}Traffic; passing\textquotedblright{}.) Say it once more.
\end{tcolorbox}
\section[\texorpdfstring{\ar{ازدحام} (izdiḥām)}{ازدحام (izdiḥām)}]{\ar{ازدحام} (izdiḥām) — congestion; a traffic jam}

\label{lesson:AR-C431-izdiham}



\begin{quote}
Before the new one: say the Arabic for a signal, then the Arabic for traffic.
\end{quote}

\subsection*{You'll want to know: \ar{ازدحام}}
\textbf{\ar{ازدحام}} — \emph{izdiḥām} — \textquotedblleft{}congestion; a traffic jam\textquotedblright{}.

\begin{quote}
\ar{تأخرنا} \ar{بسبب} \ar{الازدحام} \ar{في} \ar{الشارع}.
\end{quote}

\begin{quote}
\emph{taʾakhkharnā bi-sabab al-izdiḥām fī ash-shāriʿ.}
\end{quote}

\begin{quote}
We were late because of the congestion in the street.
\end{quote}

\begin{cousinweb}
\textbf{\ar{ازدحام}} (\emph{izdiḥām}) is built on the native Arabic root \emph{z-ḥ-m}, \textquotedblleft{}to crowd, to press\textquotedblright{} (Lane). It is the verbal noun of Form VIII, beside \emph{muzdaḥim}, \textquotedblleft{}crowded\textquotedblright{}.
\end{cousinweb}

\begin{grammarlens}[title={What you've built}]
One more word for getting around.
\end{grammarlens}

\subsection*{Your turn}
\begin{itemize}
\raggedright
  \item \emph{Say it:} \emph{izdiḥām}
  \item \emph{Say it:} \emph{izdiḥām}, once more
  \item \emph{Say it:} \emph{taʾakhkharnā bi-sabab al-izdiḥām fī ash-shāriʿ.} — We were late because of the congestion in the street.
  \item \emph{Recall:} say the Arabic for a signal, then the Arabic for traffic, then say \emph{izdiḥām} again
\end{itemize}

\begin{tcolorbox}[breakable,colback=teal!4,colframe=teal!35!black,title={Before you move on}]
What does \ar{ازدحام} mean? (\textquotedblleft{}Congestion; a traffic jam\textquotedblright{}.) Say it once more.
\end{tcolorbox}
\section[\texorpdfstring{\ar{انعطف} (inʿaṭafa)}{انعطف (inʿaṭafa)}]{\ar{انعطف} (inʿaṭafa) — to turn off; to turn (left or right)}

\label{lesson:AR-C431-inatafa}



\begin{quote}
Before the new one: say the Arabic for traffic, then the Arabic for congestion.
\end{quote}

\subsection*{You'll want to know: \ar{انعطف}}
\textbf{\ar{انعطف}} — \emph{inʿaṭafa} — \textquotedblleft{}to turn off; to turn (left or right)\textquotedblright{}.

\begin{quote}
\ar{انعطفنا} \ar{يسارا} \ar{بعد} \ar{الجسر}.
\end{quote}

\begin{quote}
\emph{inʿaṭafnā yasāran baʿda al-jisr.}
\end{quote}

\begin{quote}
We turned left after the bridge.
\end{quote}

Present: \ar{ينعطف}, \emph{yanʿaṭif}.

\begin{cousinweb}
\textbf{\ar{انعطف}} (\emph{inʿaṭafa}) is built on the native Arabic root \emph{ʿ-ṭ-f}, \textquotedblleft{}to bend, to turn\textquotedblright{} (Lane). It is Form VII.
\end{cousinweb}

\begin{grammarlens}[title={What you've built}]
One more word for getting around.
\end{grammarlens}

\subsection*{Your turn}
\begin{itemize}
\raggedright
  \item \emph{Say it:} \emph{inʿaṭafa}
  \item \emph{Say it:} \emph{inʿaṭafa}, once more
  \item \emph{Say it:} \emph{inʿaṭafnā yasāran baʿda al-jisr.} — We turned left after the bridge.
  \item \emph{Recall:} say the Arabic for traffic, then the Arabic for congestion, then say \emph{inʿaṭafa} again
\end{itemize}

\begin{tcolorbox}[breakable,colback=teal!4,colframe=teal!35!black,title={Before you move on}]
What does \ar{انعطف} mean? (\textquotedblleft{}To turn off; to turn (left or right)\textquotedblright{}.) Say it once more.
\end{tcolorbox}
